% Supplementary Information of the ERL Letter (D-107): the generality test at Great Salt Lake
% (brief section 11, PREREG Amendment 16, D-094 and addenda). Every number is a \metric{} macro.

\subsection*{Design}
The design was fixed in pre-registration Amendment~16 before any water frequency, observed extent or hindcast quantity of either candidate basin existed. Great Salt Lake (Utah, USA) was run first and Lake Urmia second; the Aral Sea was not used. The pipeline is the Caspian one in lighter form: Landsat only, \metric{params.stage2_res_m}~m counts, per-acquisition thresholds, level bins of \metric{params.bin_width_m}~m, observed extents connected to an open-water seed, a noise floor from two interleaved halves of the acquisitions, the flood-fill hindcast and the Caspian decision rule unchanged, with the lake's two arms as the sectors. The verdict rule was registered with it: the pattern ``recurs'' if at least three of the five grids mispredict in a majority of a basin's evaluable sectors, ``does not recur'' if at most one does, and is ``partial'' if two do or if the positive control mispredicts.

\emph{Arms and levels.} The north and south arms are separated by the railway causeway of 1959 and the Promontory medial line; Farmington Bay and Bear River Bay are excluded (river-fed, impounded and ungauged), and each arm is its own hindcast, with grid heights set to missing outside it. Levels are the daily means of the USGS gauges at Saltair (south arm) and Saline (north arm)\citep{usgs_nwis}, converted from the National Geodetic Vertical Datum of 1929 to the North American Vertical Datum of 1988 (NAVD~88) with VERTCON~3.0 (\metric{generality.great_salt_lake.north.ngvd29_to_navd88_m}~m at Saline, \metric{generality.great_salt_lake.south.ngvd29_to_navd88_m}~m at Saltair), computed with the NOAA coordinate conversion tool NCAT at each gauge\citep{noaa_ncat}, bridged linearly across gaps of at most 16 days. The rising limb E1 ends on \metric{generality.great_salt_lake.south.e1_end} (south arm) and \metric{generality.great_salt_lake.north.e1_end} (north arm). In all, \metric{generality.great_salt_lake.n_scenes} Landsat scenes in \metric{generality.great_salt_lake.n_acquisitions} acquisitions entered. All areas use a Lambert azimuthal equal-area projection centred on the lake.

\emph{Total variant.} The grids were brought to NAVD~88 per cell through the difference between NAVD~88 and EGM2008 heights computed with PROJ through the ITRF2014 to NAD83(2011) transformation; PROJ's direct route drops the shift between the NAD83 and WGS84 ellipsoids.

\emph{Positive control.} The pre-registered control, the USGS topobathymetric elevation model of the lake\citep{root2023_gsl_tbdem}, could not be retrieved by script, so its pre-registered substitute was used: the 1/3 arc-second elevation model of the USGS 3D Elevation Program\citep{usgs_3dep10m}, a hydro-flattened mosaic of sources\citep{archuleta2017,usgs_3dep13_tiles} that over the lake-bed band is mostly lidar acquired between September and November 2016\citep{quantumspatial2017_gsl}, evaluated only at level bins wholly above its hydro-flattened lake surfaces. It tests the audit against topography, mostly lidar, acquired after most of the band had emerged, not against the pre-emergence survey; it covers only the upper levels; and it shares the observed extents and noise floor with the grid audit.

\emph{Cost.} The Earth Engine use was \metric{generality.great_salt_lake.eecu_h_spent}~EECU-h against a cap of \metric{generality.great_salt_lake.eecu_h_cap}~EECU-h. Lake Urmia was not run within the month's budget (\metric{generality.lake_urmia.verdict}).

\subsection*{Results}
The verdict is ``\metric{generality.great_salt_lake.verdict}'' in the shape variant and ``\metric{generality.great_salt_lake.verdict_total_variant}'' in the total variant: \metric{generality.great_salt_lake.n_products_mispredict} of \metric{generality.great_salt_lake.n_products_evaluable} grids mispredict both arms, and the positive control passes (table~\ref{tab:si_gsl}). In the total variant, four grids put their flat lake surface above most observed levels, so their flood fill predicts no water there and their displacement is infinite. ETOPO 2022's emerged band is \metric{erl.text.gsl.north.etopo2022_fabdem_pct}\% FABDEM in the north arm and \metric{erl.text.gsl.south.etopo2022_fabdem_pct}\% in the south arm; its displacement is \metric{erl.text.m0.generality.great_salt_lake.etopo_2022.shape.north.median_displacement_above_2013_m}~m above the level of 1~July 2013 (\metric{generality.great_salt_lake.north.level_2013_07_01_m}~m) and \metric{erl.text.m0.generality.great_salt_lake.etopo_2022.shape.north.median_displacement_below_2013_m}~m below it in the north arm, and \metric{erl.text.m0.generality.great_salt_lake.etopo_2022.shape.south.median_displacement_above_2013_m} and \metric{erl.text.m0.generality.great_salt_lake.etopo_2022.shape.south.median_displacement_below_2013_m}~m in the south arm. It is the weakest failure: in the north arm it exceeds the floor at \metric{generality.great_salt_lake.etopo_2022.shape.north.n_exceeding_levels} of \metric{generality.great_salt_lake.etopo_2022.shape.north.n_evaluable_levels} levels, and with a floor pooled over all levels it does not mispredict that arm. Every other grid--arm pair mispredicts under both floors.

\begin{table}[!htbp]
\centering
\small
\caption{\textbf{Great Salt Lake: median shoreline displacement after the best-fitting vertical constant (m).} Median over each arm's evaluable levels of the per-level median distance between predicted and observed shorelines, and the number of levels at which it exceeds the \metric{params.err_pct_hi}th percentile of the noise floor. Noise floor: median over levels of that percentile, \metric{erl.text.m0.generality.great_salt_lake.north.noise_floor_p95_median_m}~m (north arm) and \metric{erl.text.m0.generality.great_salt_lake.south.noise_floor_p95_median_m}~m (south arm); for the 3DEP control, at its levels only, \metric{erl.text.m0.generality.great_salt_lake.control.north.median_noise_p95_m} and \metric{erl.text.m0.generality.great_salt_lake.control.south.median_noise_p95_m}~m.}
\label{tab:si_gsl}
\begin{tabular}{lrrrrl}
\hline
Surface & North (m) & Levels over floor & South (m) & Levels over floor & Arms mispredicted, shape / total \\
\hline
GEBCO\_2026 & \metric{erl.text.m0.generality.great_salt_lake.gebco_2026.shape.north.median_displacement_m} & \metric{generality.great_salt_lake.gebco_2026.shape.north.n_exceeding_levels}/\metric{generality.great_salt_lake.gebco_2026.shape.north.n_evaluable_levels} & \metric{erl.text.m0.generality.great_salt_lake.gebco_2026.shape.south.median_displacement_m} & \metric{generality.great_salt_lake.gebco_2026.shape.south.n_exceeding_levels}/\metric{generality.great_salt_lake.gebco_2026.shape.south.n_evaluable_levels} & \metric{generality.great_salt_lake.gebco_2026.shape.n_sectors_mispredict} / \metric{generality.great_salt_lake.gebco_2026.total.n_sectors_mispredict} \\
GEBCO\_2022 & \metric{erl.text.m0.generality.great_salt_lake.gebco_2022.shape.north.median_displacement_m} & \metric{generality.great_salt_lake.gebco_2022.shape.north.n_exceeding_levels}/\metric{generality.great_salt_lake.gebco_2022.shape.north.n_evaluable_levels} & \metric{erl.text.m0.generality.great_salt_lake.gebco_2022.shape.south.median_displacement_m} & \metric{generality.great_salt_lake.gebco_2022.shape.south.n_exceeding_levels}/\metric{generality.great_salt_lake.gebco_2022.shape.south.n_evaluable_levels} & \metric{generality.great_salt_lake.gebco_2022.shape.n_sectors_mispredict} / \metric{generality.great_salt_lake.gebco_2022.total.n_sectors_mispredict} \\
ETOPO 2022 & \metric{erl.text.m0.generality.great_salt_lake.etopo_2022.shape.north.median_displacement_m} & \metric{generality.great_salt_lake.etopo_2022.shape.north.n_exceeding_levels}/\metric{generality.great_salt_lake.etopo_2022.shape.north.n_evaluable_levels} & \metric{erl.text.m0.generality.great_salt_lake.etopo_2022.shape.south.median_displacement_m} & \metric{generality.great_salt_lake.etopo_2022.shape.south.n_exceeding_levels}/\metric{generality.great_salt_lake.etopo_2022.shape.south.n_evaluable_levels} & \metric{generality.great_salt_lake.etopo_2022.shape.n_sectors_mispredict} / \metric{generality.great_salt_lake.etopo_2022.total.n_sectors_mispredict} \\
ETOPO1 & \metric{erl.text.m0.generality.great_salt_lake.etopo1.shape.north.median_displacement_m} & \metric{generality.great_salt_lake.etopo1.shape.north.n_exceeding_levels}/\metric{generality.great_salt_lake.etopo1.shape.north.n_evaluable_levels} & \metric{erl.text.m0.generality.great_salt_lake.etopo1.shape.south.median_displacement_m} & \metric{generality.great_salt_lake.etopo1.shape.south.n_exceeding_levels}/\metric{generality.great_salt_lake.etopo1.shape.south.n_evaluable_levels} & \metric{generality.great_salt_lake.etopo1.shape.n_sectors_mispredict} / \metric{generality.great_salt_lake.etopo1.total.n_sectors_mispredict} \\
SRTM15+ V2.7 & \metric{erl.text.m0.generality.great_salt_lake.srtm15p_v2.7.shape.north.median_displacement_m} & \metric{generality.great_salt_lake.srtm15p_v2.7.shape.north.n_exceeding_levels}/\metric{generality.great_salt_lake.srtm15p_v2.7.shape.north.n_evaluable_levels} & \metric{erl.text.m0.generality.great_salt_lake.srtm15p_v2.7.shape.south.median_displacement_m} & \metric{generality.great_salt_lake.srtm15p_v2.7.shape.south.n_exceeding_levels}/\metric{generality.great_salt_lake.srtm15p_v2.7.shape.south.n_evaluable_levels} & \metric{generality.great_salt_lake.srtm15p_v2.7.shape.n_sectors_mispredict} / \metric{generality.great_salt_lake.srtm15p_v2.7.total.n_sectors_mispredict} \\
3DEP (control) & \metric{erl.text.m0.generality.great_salt_lake.control.north.median_displacement_m} & \metric{generality.great_salt_lake.control.north.n_exceeding_levels}/\metric{generality.great_salt_lake.control.north.n_evaluable_levels} & \metric{erl.text.m0.generality.great_salt_lake.control.south.median_displacement_m} & \metric{generality.great_salt_lake.control.south.n_exceeding_levels}/\metric{generality.great_salt_lake.control.south.n_evaluable_levels} & \metric{generality.great_salt_lake.positive_control} \\
\hline
\end{tabular}
\end{table}
